% MP3_Graph_Theory_Network_Models.tex
\documentclass[article, 11pt]{article}

%%%%%%%%%%%%%%%%%%%%%%%%%%%%%%%%%%%%%%%%%%%%%%%%%%%
% Preamble:
%%%%%%%%%%%%%%%%%%%%%%%%%%%%%%%%%%%%%%%%%%%%%%%%%%%
\usepackage{fullpage}
\usepackage{amsfonts}
\usepackage{amsmath}
\usepackage{amsthm}
\usepackage{graphicx}
\usepackage{color}
\usepackage{palatino, url, multicol}
\usepackage{enumerate}
\usepackage{ulem}
\usepackage{hyperref}
\usepackage{verbatim}

\thispagestyle{empty}

%%%%%%%%%%%%%%%%%%%%%%%%%%%%%%%%%%%%%%%%%%%%%%%%%%%%
% Start Document:
%%%%%%%%%%%%%%%%%%%%%%%%%%%%%%%%%%%%%%%%%%%%%%%%%%%%
\begin{document}

% Project Title:
\begin{center}
\boxed{\LARGE{\textbf{\textsc{MINI-PROJECT 3}}}}\\
\Huge{\textbf{\textsc{Graph Theory Network Models}}}\\
\end{center}

\vspace{.6cm}

% Brief Summary:
\noindent \boxed{\textbf{\textsc{Brief Summary}}}

\vspace{-.35cm}
\noindent \hrulefill

This mini-project uses university enrollment data to build and analyze a network of academic departments connected by students completing majors and minors. For each year from 2006 to 2024, you will construct a weighted, undirected department network where edge weights measure how strongly two departments are connected by students with layers and an interaction parameter: 
\[
A(\alpha) = A_{\text{mm}} + \alpha A_{\text{mn}} + \alpha^2 A_{\text{nn}},
\] 

You will investigate two main themes: (1) Community structure over time -- detect groups of departments that are more connected internally than externally, and track how community structure changes over time using modularity; and (2) Department centrality over time -- compute centrality measures each year and interpret which departments act as ``hubs'' or ``bridges'' in the curriculum ecosystem. A key goal is connecting graph measures (centrality, modularity, communities) to a real institutional setting, and communicating trends clearly using figures.

\vspace{.35cm}


% Goals:
\noindent \boxed{\textbf{\textsc{Goals}}}

\vspace{-.35cm}
\noindent \hrulefill

\begin{itemize}
\item Build a weighted department network for each year from 2006 to 2024 using Python.
\item Compute and interpret {centrality measures over time}, such as: strength (weighted degree), closeness centrality, betweenness centrality, eigenvector centrality. 
\item Detect {community structure} each year using a modularity-based method, and track: the number of communities over time, modularity over time, notable shifts in membership .
\item Produce clear, well-labeled visualizations that communicate network structure and trends.
\end{itemize}

\vspace{.35cm}

% LaTeX Deliverables:
\noindent \boxed{\textbf{\textsc{\LaTeX\, Deliverables:}}}

\vspace{-.26cm}
\noindent \hrulefill

\begin{itemize}
\item A short \LaTeX\, report that is clearly written and well-organized.
\item \text{Include at least three figures with captions}:
  \begin{enumerate}
    \item \text{Graph figure:} a visualization of the department network for at least one selected year (or multiple years side-by-side). Use edge width/opacity to represent weight and a consistent layout when comparing years.
    \item \text{Centrality figure:} a plot showing at least one centrality measure over time (2006--2024). Examples: top 5 departments by eigenvector centrality each year; or a line plot for a small set of departments you select.
    \item \text{Communities/modularity figure:} a plot of \text{modularity over time}. If possible, also show number of communities over time or a community membership heatmap/diagram for selected years.
  \end{enumerate}
\item Provide interpretation: explain \text{why} certain departments appear central, and what community structure suggests about clusters of programs.
\end{itemize}

\vspace{.35cm}

% Python Deliverables:
\noindent \boxed{\textbf{\textsc{Python Deliverables:}}}

\vspace{-.35cm}
\noindent \hrulefill

\begin{itemize}
\item A Python script or notebook that:
  \begin{itemize}
  \item loads the data for 2006--2024,
  \item constructs the network(s) each year,
  \item computes centrality metrics each year,
  \item detects communities each year and computes modularity,
  \item generates and saves the required figures.
  \end{itemize}
\item Use \texttt{networkx} for graph construction and centralities. For community detection, you may use:
  \begin{itemize}
    \item \texttt{networkx.algorithms.community} methods, or
    \item another standard Python community detection package.
  \end{itemize}
\item All plots must have titles, axis labels, legends, and readable formatting.
\end{itemize}

\vspace{.35cm}

% Suggested Analysis Questions (optional):
\noindent \boxed{\textbf{\textsc{Suggested Analysis Questions}}}

\vspace{-.35cm}
\noindent \hrulefill

\begin{itemize}
\item Which departments are consistently central across years? Which departments rise/fall sharply?
\item Do different centralities (strength vs.\ betweenness vs.\ eigenvector) tell different stories? Give at least one example.
\item When modularity changes, what changes in the network might explain it (new bridges, merging/splitting clusters, etc.)?
\item Are communities stable over time, or do departments frequently change communities?
\end{itemize}

\vspace{.35cm}

% AI Disclosure:
\noindent \boxed{\textbf{\textsc{AI Disclosure:}}}

\vspace{-.35cm}
\noindent \hrulefill

If you use AI in any way, you must include an Appendix in the write-up that notes the LLM used and provides a list of the prompts you used. You may use AI for suggestions related to Python code, \LaTeX\, syntax, or other general ideas, but \textit{the report should not be copy-and-pasted from AI.}

\pagebreak

\end{document}